\documentclass[14pt,a4paper]{extarticle}
\usepackage[left=25mm,right=15mm,top=20mm,bottom=20mm,headheight=18pt]{geometry}
\usepackage{fontspec}
\usepackage{polyglossia}
\setmainlanguage{ukrainian}
\setotherlanguage{english}
\setmainfont{Times New Roman}
\newfontfamily\cyrillicfont{Times New Roman}
\setmonofont{Courier New}
\newfontfamily\cyrillicfonttt{Courier New}
\usepackage{setspace}
\onehalfspacing
\usepackage{indentfirst}
\setlength{\parindent}{1.25cm}
\setlength{\parskip}{0pt}
\usepackage{amsmath}
\usepackage{graphicx}
\usepackage{caption}
\DeclareCaptionLabelSeparator{dstu}{ — }
\captionsetup{font=normalsize,labelsep=dstu}
\captionsetup[figure]{name=Рисунок,justification=centering}
\captionsetup[table]{name=Таблиця,position=top,justification=raggedright,
  singlelinecheck=false,margin={1.25cm,0cm}}
\usepackage{titlesec}
\titleformat{\section}{\centering\normalfont\bfseries}{\thesection}{1em}{}
\titleformat{\subsection}{\normalfont\bfseries}{\thesubsection}{1em}{}
\titlespacing*{\section}{0pt}{2\baselineskip}{2\baselineskip}
\titlespacing*{\subsection}{1.25cm}{2\baselineskip}{2\baselineskip}
\usepackage{fancyhdr}
\pagestyle{fancy}
\fancyhf{}
\fancyhead[R]{\thepage}
\renewcommand{\headrulewidth}{0pt}
\usepackage[hidelinks]{hyperref}
\renewcommand{\labelitemi}{—}
\renewcommand{\refname}{ПЕРЕЛІК ДЖЕРЕЛ ПОСИЛАННЯ}
\emergencystretch=2em

\graphicspath{{./pictures/}}
\newcommand{\ministry}{МІНІСТЕРСТВО ОСВІТИ І НАУКИ УКРАЇНИ}
\newcommand{\university}{НАЦІОНАЛЬНИЙ УНІВЕРСИТЕТ «ЛЬВІВСЬКА ПОЛІТЕХНІКА»}
\newcommand{\institute}{Інститут комп'ютерних наук та інформаційних технологій}
\newcommand{\department}{систем штучного інтелекту}
\newcommand{\subject}{[Назва дисципліни]}
\newcommand{\lecturer}{{[Посада лектора]}\\{}[Прізвище та ініціали]}
\newcommand{\teacher}{{[Посада викладача]}\\{}[Прізвище та ініціали]}
\newcommand{\mygroup}{[Шифр групи]}
\newcommand{\studentname}{{[Ваше прізвище та ініціали]}}
\newcommand{\studentdescription}{студент групи}
\newcommand{\studentrole}{Виконав}
\newcommand{\teacherrole}{Прийняв}
\newcommand{\lab}{[Номер]}
\newcommand{\variant}{[Номер]}
\newcommand{\theme}{[Тема лабораторної роботи]}
\newcommand{\purpose}{[Мета лабораторної роботи]}
\newcommand{\city}{Львів}
\newcommand{\reportyear}{\the\year}

\begin{document}
\newgeometry{left=20mm,right=20mm,top=15mm,bottom=15mm}
\begin{titlepage}
    \singlespacing
    \thispagestyle{empty}
    \begin{center}
        \textbf{\ministry\\\university}
    \end{center}
    \vspace{10pt}
    \begin{center}
        \institute\\[6pt]
        Кафедра \department
    \end{center}
    \begin{center}
        \includegraphics[width=60mm]{university-logo.png}
    \end{center}
    \vspace{10pt}
    \begin{center}
        \textbf{ЗВІТ}\\[10pt]
        до лабораторної роботи №~\lab\\
        \textbf{на тему}: «\textit{\theme}»\\
        \textbf{з дисципліни}: «\subject»
    \end{center}
    \vspace{16pt}
    \begin{flushright}
        \textbf{Лектор}:\\
        \lecturer\\[16pt]
        \textbf{\studentrole}:\\
        \studentdescription{} \mygroup\\
        \studentname\\[16pt]
        \textbf{\teacherrole}:\\
        \teacher\\[16pt]
        «\rule{1cm}{0.15mm}» \rule{1.5cm}{0.15mm} \reportyear{} р.\\
        $\sum = {}$\rule{1cm}{0.15mm}\ldots\ldots\ldots
    \end{flushright}
    \vfill
    \begin{center}
        \textbf{\city{} — \reportyear}
    \end{center}
\end{titlepage}
\restoregeometry
\setcounter{page}{2}

\section*{ТЕМА РОБОТИ}
[Тема з методичних вказівок.]

\section*{МЕТА РОБОТИ}
[Мета з методичних вказівок або погоджене формулювання.]

\section{ЗАВДАННЯ}
[Умова, варіант, вхідні дані та очікуваний результат.]

\section{ХІД РОБОТИ ТА РЕЗУЛЬТАТИ}
[Опис фактично виконаних кроків, код, вимірювання та їх пояснення.]

% Додайте лише фактичні результати. Приклади конструкцій:
% На рисунку~\ref{fig:result} показано результат виконання.
% \begin{figure}[htbp]
%   \centering
%   \includegraphics[width=0.85\linewidth]{images/result.png}
%   \caption{Результат виконання}\label{fig:result}
% \end{figure}
%
% Дані наведено в таблиці~\ref{tab:results}.
% \begin{table}[htbp]
%   \caption{Результати вимірювань}\label{tab:results}
%   \centering
%   \begin{tabular}{ll}
%     Параметр & Значення \\
%     Параметр 1 & Фактичне значення
%   \end{tabular}
% \end{table}
%
% Середнє значення обчислено за формулою~\eqref{eq:mean}.
% \begin{equation}\label{eq:mean}
%   \bar{x}=\frac{1}{n}\sum_{i=1}^{n}x_i
% \end{equation}
% \noindent де $x_i$ — результат вимірювання; $n$ — кількість вимірювань.

\section*{ВИСНОВКИ}
[Що виконано, які результати отримано та чи досягнуто мети.]

% Якщо використано джерела, додайте \cite{key} у тексті та
% thebibliography з \bibitem{key} у порядку першого згадування.
\end{document}
